\documentclass[11pt]{article}
\usepackage[margin=1in]{geometry}
\usepackage{amsmath,amssymb,amsthm}
\usepackage[colorlinks=true,urlcolor=blue,linkcolor=blue]{hyperref}
\newtheorem{theorem}{Theorem}
\newtheorem{lemma}[theorem]{Lemma}
\newtheorem{proposition}[theorem]{Proposition}
\newcommand{\R}{\mathbb R}
\newcommand{\norm}[1]{\left\lVert #1\right\rVert}
\newcommand{\range}{\operatorname{range}}
\newcommand{\rank}{\operatorname{rank}}
\newcommand{\tO}{\widetilde O}
\title{Clustered tails: a conditional product-approximation theorem and why it is not a new Krylov mechanism}
\author{Research audit}
\date{November 14, 2026}
\begin{document}
\maketitle

\paragraph{Status.}
The conditional theorem below is proved in the exact-real arithmetic model,
using the established constant-accuracy randomized block Krylov theorem.
It is \emph{not} the desired arbitrary-input theorem, and it is not evidence of
a new state-of-the-art algorithm. Its acceleration is explained by the
classical affine invariance of Krylov spaces and clustered-spectrum
convergence. In particular, an apparent need for a new product-aware
projector primitive disappears once iterates are kept in the range of the
left factor. No finite-precision complexity theorem is claimed here.

\section{Setting and conditional theorem}
Let \(A\in\R^{n\times m}\), \(B\in\R^{m\times p}\),
\(M=AB\), \(H=MM^\top\), \(I=m(n+p)\), and
\(a=\rank A\). Write \(\lambda_i=\sigma_i(M)^2\).
Fix \(1\le r<a\) and \(0<\theta<1\), and put
\(t=\lambda_{r+1}>0\). The additional assumption is
\begin{equation}\label{eq:cluster}
 \lambda_a\ge(1-\rho)t,\qquad 0\le\rho\le\theta.
\end{equation}
Thus \(M\) has full row rank on \(\range A\), and its squared
tail singular values lie in \([(1-\rho)t,t]\). The leading
\(r\) values are arbitrary. This assumption may fail badly for a
general product, for instance when \(B\) annihilates part of
\(\range A^\top\).

\begin{theorem}[Conditional clustered-tail result]\label{thm:cluster}
Assume \eqref{eq:cluster}. Given constant-distortion input embeddings
with their usual \(\tO(I+m^{\omega+\zeta})\) setup, there is a
randomized exact-real arithmetic algorithm producing rank-at-most-\(r\)
factors and their explicit product \(C\) such that
\[
 \Pr\{\norm{M-C}_2\le(1+\theta)\sigma_{r+1}(M)\}\ge1-\delta.
\]
Its construction costs
\(\tO(Ir+m^{\omega+\zeta})\), and expansion costs \(O(npr)\).
The threshold \(t\) need not be supplied.
\end{theorem}

The proof below first gives a version with an approximate projector. That
version isolates a useful exact selfadjoint polynomial and avoids silently
assuming an available orthonormal basis for \(\range A\). A subsequent
simplification eliminates the projector actions and shows directly why the
mechanism is classical Krylov acceleration.

\section{Estimating a centering threshold with \(r+1\) probes}
Let \(Z\in\R^{n\times(r+1)}\) have full column rank and lie in
\(\range A\). It can be generated as \(Z=A\Omega\) for a
Gaussian \(\Omega\), with probability one in exact arithmetic; a
preconditioned choice is preferable for controlling its conditioning.
Compute
\[
 G=Z^\top Z,\qquad K=(M^\top Z)^\top(M^\top Z),\qquad
 t_0=\lambda_{\min}(K,G).
\]
Here the last expression is the smallest generalized eigenvalue, equivalently
the minimum Rayleigh quotient of \(H\) on \(\range Z\).

\begin{lemma}\label{lem:threshold}
For every such full-rank \(Z\),
\[
 (1-\rho)t\le t_0\le t.
\]
The upper bound \(t_0\le\lambda_{r+1}(H)\) holds without the
cluster assumption.
\end{lemma}
\begin{proof}
On \(\range A\), the smallest eigenvalue of \(H\) is at
least \((1-\rho)t\), proving the lower bound. The
\((r+1)\)-dimensional space \(\range Z\) intersects the orthogonal
complement of the first \(r\) eigenvectors nontrivially. A vector in
that intersection has Rayleigh quotient at most \(t\).
\end{proof}

Generating \(Z\), applying \(M^\top\), and forming the two small
Gram matrices cost \(O(I(r+1)+(n+p)(r+1)^2)\). Since \(r<m\),
this is \(O(Ir)\) up to constants. Small-core generalized eigensolving
costs at most \(\tO(m^{\omega+\zeta})\). Exact arithmetic is
material: a conventional floating-point computation of a very small
\(t_0\) in the presence of arbitrarily large head eigenvalues needs
an explicit precision analysis.

\section{An initially proposed approximate-projector construction}
Let \(P=A^\top A\), and let \(\widetilde P\) satisfy
\((1-\eta)P\preceq\widetilde P\preceq(1+\eta)P\), where
\(0<\eta<1/2\) is fixed. Let \(\Pi_A\) be the orthogonal
projector onto \(\range A\), and define
\[
 J=A\widetilde P^\dagger A^\top,\qquad c=1-\eta^2.
\]
The nonzero eigenvalues of \(J\) lie in
\([1/(1+\eta),1/(1-\eta)]\); the kernels of \(P\) and
\(\widetilde P\) agree.

For an even integer \(j\), define the fixed selfadjoint polynomial
\begin{equation}\label{eq:projector}
 \widehat\Pi_A=I-(I-cJ)^j.
\end{equation}
Although an identity appears here, the polynomial vanishes on
\((\range A)^\perp\).
\begin{lemma}\label{lem:projector}
For even \(j\),
\[
 0\preceq\widehat\Pi_A\preceq\Pi_A,
 \qquad \norm{\widehat\Pi_A-\Pi_A}_2\le\eta^j.
\]
Its action on a block of \(s\) vectors costs \(O(Isj)\) after
the small-core inverse has been formed.
\end{lemma}
\begin{proof}
On \(\range A\), every eigenvalue of \(I-cJ\) lies in
\([-\eta,\eta]\). Even powers are positive semidefinite with norm
at most \(\eta^j\). On the complementary space the same matrix
equals the identity, so \eqref{eq:projector} is zero there. Each
application of \(J\) uses \(A^\top\), a dense small-core inverse,
and \(A\); the \(O(m^2s)\) middle cost is absorbed by \(O(Is)\).
\end{proof}

Choose even \(j=O(\log(1/\theta))\) so that
\(\xi=\eta^j\le\theta^2/2\), and set
\[
 L=H-t_0\Pi_A,\qquad
 \widehat L=H-t_0\widehat\Pi_A.
\]
On the tail of \(H\) inside \(\range A\), every eigenvalue of
\(L\) has magnitude at most \(\rho t\). Outside \(\range A\),
\(L\) is zero. Thus \(\sigma_{r+1}(L)\le\rho t\), and Weyl's
singular-value inequality gives
\[
 \sigma_{r+1}(\widehat L)\le(\rho+\xi)t.
\]
Apply the known constant-relative-error randomized block Krylov theorem to
the symmetric (possibly indefinite) operator \(\widehat L\). With
polylogarithmically many block applications, it returns orthonormal
\(U\in\R^{n\times r}\) satisfying
\[
 \norm{(I-UU^\top)\widehat L}_2\le2(\rho+\xi)t
\]
with the assigned failure probability. The theorem is applied to singular
values, so indefiniteness causes no logical problem.

Return \(C=UU^\top M\). Since
\(H=t_0\widehat\Pi_A+\widehat L\) and
\(0\preceq\widehat\Pi_A\preceq I\),
\begin{align*}
 \norm{M-C}_2^2
 &=\norm{(I-UU^\top)H(I-UU^\top)}_2\\
 &\le t_0+\norm{(I-UU^\top)\widehat L(I-UU^\top)}_2\\
 &\le t+2(\rho+\xi)t
 \le(1+\theta)^2t.
\end{align*}
This proves the claimed error bound. Use the factors
\(X=U\), \(Y=M^\top U\); obtaining \(Y\) costs \(O(Ir)\).

There are \(\tO(r)\) Krylov basis vectors and input passes. QR costs
\(\tO((n+p)r^2)\), absorbed by \(\tO(Ir)\) because \(r\le m\).
If a conventional block implementation would grow past dimension \(m\),
cap it at the invariant-space dimension and use exact processing there;
alternatively the theorem's polylogarithmic factors allow the core algebra
bound directly. Dense small projected algebra costs
\(\tO(m^{\omega+\zeta})\), using exponent slack to absorb logarithms.
The even-\(j\) projector construction is fixed before Krylov starts,
so there is no substitution of adaptively inaccurate nonlinear actions for
a fixed symmetric operator.

Allocate failure probability between the input embedding and the Krylov
start. Threshold generation fails only on a measure-zero event in the
exact Gaussian model. Expansion costs \(O(npr)\) in addition to every
construction cost above.

\section{Why the projector is unnecessary}
On \(\range A\), \(\Pi_A\) is the identity. Choose the initial
block within this space and preserve that invariant space in exact
arithmetic. Then use
\[
 (H-t_0I)v=H v-t_0v,\qquad v\in\range A,
\]
instead of computing a projector. The unwanted ambient eigenvalue
\(-t_0\) on \((\range A)^\perp\) never enters the iterates.

A well-conditioned random start can be generated from a constant sketch.
On the \(a\)-dimensional support of \(\widetilde P\), let
\(G_0\) be standard Gaussian and set
\(Z_0=A\widetilde P^{\dagger/2}G_0\). Its covariance restricted
to \(\range A\) has eigenvalues in a fixed constant interval.
The usual Gaussian starting-matrix estimates therefore lose only a
constant covariance condition number. More explicitly, the leading
\(r\)-row block has covariance bounded below by a constant, while
the full norm is bounded above by a constant times the full Gaussian norm;
independence of those two blocks is unnecessary for a union bound.

For this implementation, the preceding error proof uses
\(\sigma_{r+1}(L)\le\rho t\) directly, and gives the slightly
stronger bound
\[
 \norm{M-C}_2^2\le t+2\rho t\le(1+\theta)^2t.
\]
The projector polynomial and its \(\log(1/\theta)\) inner factor
are thus optional. They do not supply the source of the acceleration.

\paragraph{Affine invariance.}
For every fixed start \(Z_0\),
\[
 \mathcal K_q(H-t_0I,Z_0)=\mathcal K_q(H,Z_0),
\]
by the binomial formula and its inverse. The singular-value Krylov space
for \(L\), containing powers \(L,L^3,\ldots\), is contained in an
ordinary \(H\)-Krylov space of comparable degree. Consequently the
conditional theorem uses no information unavailable to classical
polynomial Krylov iteration. It changes the convergence analysis and the
selection of the final rank-\(r\) subspace within that space.

In floating point, a naive ambient-space implementation may introduce
components outside \(\range A\), where the shift has eigenvalue
\(-t_0\). One can retain coefficient representations and regenerate
tall vectors from \(A\), or use the fixed approximate-projector
version. Neither observation alone proves a stable bit complexity bound.
Conditioning of \(A\), scale separation \(\norm M^2/t\), and
loss of orthogonality have to be tracked explicitly.

\section{Hard-regime accounting and novelty audit}
In the supplied regime
\(n=p=N\), \(m=N^{8/9}\), \(r=1\), \(\theta=1/m\),
the additional hypothesis is
\[
 \sigma_m(M)^2\ge(1-1/m)\sigma_2(M)^2.
\]
Under that hypothesis, construction is
\(\tO(N^{17/9}+N^{2+(8/9)\zeta})\), and explicit expansion is
\(O(N^2)\). This equals the target exponent for this input class.
It improves the displayed \emph{worst-case upper bounds} for unadapted
Krylov and Gram routes on that class. It does not prove an improvement
over the best previously known spectrum-adaptive Krylov method on the
same class.

Martinsson--Tropp's survey explicitly records affine invariance of Krylov
spaces and the classical role of blocks for clustered eigenvalues
(Sections 6.4--6.5). Musco--Musco give gap-free low-rank guarantees and
spectrum-dependent improvements (Theorems 10 and 13). The theorem here
is a short application of those tools after centering on the invariant
space. The threshold estimator of Lemma~\ref{lem:threshold} is a useful
explicit min--max observation, but no novelty claim is supported by the
search performed so far. It should not be advertised as a new
state-of-the-art AMM result.

\paragraph{What remains open.}
The arbitrary-input target must handle a tail spread across a constant
interval as well as a narrow cluster. A signed shift then leaves a
large negative tail, and constant-accuracy approximation of the shifted
operator no longer implies \(O(\theta t)\) residual. The conditional
proof supplies no universal threshold-local preconditioner, no cheap
interior resolvent, and no arithmetic lower bound against such methods.

\begin{thebibliography}{9}
\bibitem{mm}
C. Musco and C. Musco.
\emph{Randomized Block Krylov Methods for Stronger and Faster Approximate
Singular Value Decomposition} (2015), especially Theorems 10 and 13.
\url{https://arxiv.org/abs/1504.05477}.
\bibitem{mt}
P.-G. Martinsson and J. Tropp.
\emph{Randomized Numerical Linear Algebra: Foundations and Algorithms}.
Acta Numerica 29 (2020), Sections 6.4--6.5.
\url{https://arxiv.org/abs/2002.01387}.
\bibitem{demmel}
J. Demmel, I. Dumitriu, and O. Holtz.
\emph{Fast Linear Algebra Is Stable} (2007).
\url{https://arxiv.org/abs/math/0612264}.
\end{thebibliography}
\end{document}
